\documentclass[11pt,a4paper]{article}
\usepackage[margin=2.25cm]{geometry}
\usepackage{amsmath,amssymb,bm,mathtools}
\usepackage{booktabs}
\usepackage{enumitem}
\usepackage{xcolor}
\usepackage[hidelinks]{hyperref}
\hypersetup{pdftitle={Analytical Derivations for Digital Quantum Simulation of a Two-Qubit Transverse-Field Ising Model},pdfauthor={Yulan Fan},pdfsubject={Exact and Trotterized dynamics of a two-qubit transverse-field Ising model}}
\usepackage{microtype}
\usepackage{fancyhdr}
\usepackage{lmodern}
\usepackage[T1]{fontenc}

\definecolor{accent}{RGB}{30,70,120}
\newcommand{\ket}[1]{\left|#1\right\rangle}
\newcommand{\bra}[1]{\left\langle#1\right|}
\newcommand{\braket}[2]{\left\langle#1\middle|#2\right\rangle}
\newcommand{\op}[1]{\widehat{#1}}
\newcommand{\I}{\mathbb{I}}
\newcommand{\dd}{\mathrm{d}}
\newcommand{\e}{\mathrm{e}}
\newcommand{\ii}{\mathrm{i}}
\newcommand{\Order}{\mathcal{O}}

\pagestyle{fancy}
\fancyhf{}
\lhead{Two-Qubit Ising Simulation}
\rhead{Analytical Derivations}
\cfoot{\thepage}
\setlength{\headheight}{14pt}

\title{\textbf{Analytical Derivations for Digital Quantum Simulation\\of a Two-Qubit Transverse-Field Ising Model}}
\author{Yulan Fan}
\date{July 2026}

\begin{document}
\maketitle

\begin{abstract}
This note collects the hand calculations underlying a three-notebook implementation of exact and Trotterized real-time dynamics for a two-qubit transverse-field Ising model. The derivations proceed from basis states and tensor-product operators to the exact spectrum and time evolution, and then to first- and second-order product formulas implemented with $RX$ and $RZZ$ gates. The emphasis is on transparent operator algebra, independent validation, and the distinction between state error, infidelity, observable error, and circuit cost.
\end{abstract}

\tableofcontents
\newpage

\section{Model and conventions}

We study the two-qubit Hamiltonian
\begin{equation}
H=-J Z_0 Z_1-h(X_0+X_1),
\label{eq:H}
\end{equation}
where $J$ is the Ising coupling and $h$ is the transverse-field strength. We use units with $\hbar=1$ and the ordered computational basis
\begin{equation}
\bigl(\ket{00},\ket{01},\ket{10},\ket{11}\bigr).
\end{equation}
The one-qubit basis vectors and Pauli matrices are
\begin{equation}
\ket{0}=\begin{pmatrix}1\\0\end{pmatrix},\qquad
\ket{1}=\begin{pmatrix}0\\1\end{pmatrix},
\end{equation}
\begin{equation}
X=\begin{pmatrix}0&1\\1&0\end{pmatrix},\qquad
Y=\begin{pmatrix}0&-\ii\\ \ii&0\end{pmatrix},\qquad
Z=\begin{pmatrix}1&0\\0&-1\end{pmatrix}.
\end{equation}
The two-qubit basis vectors are
\begin{equation}
\ket{00}=\begin{pmatrix}1\\0\\0\\0\end{pmatrix},\quad
\ket{01}=\begin{pmatrix}0\\1\\0\\0\end{pmatrix},\quad
\ket{10}=\begin{pmatrix}0\\0\\1\\0\end{pmatrix},\quad
\ket{11}=\begin{pmatrix}0\\0\\0\\1\end{pmatrix}.
\end{equation}

\section{Two-qubit operators and Hamiltonian matrix}

The local operators are
\begin{equation}
X_0=X\otimes \I,\qquad X_1=\I\otimes X,
\end{equation}
\begin{equation}
Z_0=Z\otimes \I,\qquad Z_1=\I\otimes Z,
\end{equation}
and
\begin{equation}
Z_0Z_1=Z\otimes Z.
\end{equation}
Their explicit matrices are
\begin{equation}
X_0=
\begin{pmatrix}
0&0&1&0\\
0&0&0&1\\
1&0&0&0\\
0&1&0&0
\end{pmatrix},\qquad
X_1=
\begin{pmatrix}
0&1&0&0\\
1&0&0&0\\
0&0&0&1\\
0&0&1&0
\end{pmatrix},
\end{equation}
\begin{equation}
Z_0Z_1=\operatorname{diag}(1,-1,-1,1).
\end{equation}
A matrix can be reconstructed from its action on an ordered basis: the $j$th column is the vector obtained by applying the operator to the $j$th basis state. For example,
\begin{equation}
X_0\ket{00}=\ket{10},\quad
X_0\ket{01}=\ket{11},\quad
X_0\ket{10}=\ket{00},\quad
X_0\ket{11}=\ket{01},
\end{equation}
which immediately gives the columns of $X_0$.

Substituting these operators into Eq.~\eqref{eq:H} gives
\begin{equation}
H=
\begin{pmatrix}
-J&-h&-h&0\\
-h&J&0&-h\\
-h&0&J&-h\\
0&-h&-h&-J
\end{pmatrix}.
\label{eq:Hmatrix}
\end{equation}
The matrix is Hermitian, as required for a physical Hamiltonian.

\section{Symmetry-adapted basis and exact spectrum}

Introduce the symmetric and antisymmetric combinations
\begin{equation}
\ket{\phi_\pm}=\frac{\ket{00}\pm\ket{11}}{\sqrt{2}},\qquad
\ket{\psi_\pm}=\frac{\ket{01}\pm\ket{10}}{\sqrt{2}}.
\end{equation}
The antisymmetric states are eigenstates of the full Hamiltonian. Direct application gives
\begin{equation}
H\ket{\phi_-}=-J\ket{\phi_-},
\end{equation}
\begin{equation}
H\ket{\psi_-}=+J\ket{\psi_-}.
\end{equation}
The remaining symmetric subspace is spanned by $\{\ket{\phi_+},\ket{\psi_+}\}$. In this basis,
\begin{equation}
H_{\mathrm{sym}}=
\begin{pmatrix}
-J&-2h\\
-2h&J
\end{pmatrix}.
\end{equation}
Define
\begin{equation}
\Omega=\sqrt{J^2+4h^2}.
\end{equation}
Since
\begin{equation}
H_{\mathrm{sym}}^2=\Omega^2\I,
\end{equation}
the two eigenvalues of this block are $\pm\Omega$. The full spectrum is therefore
\begin{equation}
\boxed{E\in\{-\Omega,-J,J,\Omega\}}.
\end{equation}

\section{Exact time evolution from \texorpdfstring{$\ket{00}$}{|00>}}

The initial state decomposes as
\begin{equation}
\ket{00}=\frac{\ket{\phi_+}+\ket{\phi_-}}{\sqrt{2}}.
\end{equation}
The antisymmetric component evolves trivially:
\begin{equation}
\e^{-\ii Ht}\ket{\phi_-}=\e^{\ii Jt}\ket{\phi_-}.
\end{equation}
For the symmetric block, the identity $H_{\mathrm{sym}}^2=\Omega^2\I$ gives
\begin{equation}
\e^{-\ii H_{\mathrm{sym}}t}
=\cos(\Omega t)\I-\ii\frac{\sin(\Omega t)}{\Omega}H_{\mathrm{sym}}.
\end{equation}
Acting on $\ket{\phi_+}$,
\begin{equation}
\e^{-\ii H_{\mathrm{sym}}t}\ket{\phi_+}
=\left[\cos(\Omega t)+\ii\frac{J}{\Omega}\sin(\Omega t)\right]\ket{\phi_+}
+\ii\frac{2h}{\Omega}\sin(\Omega t)\ket{\psi_+}.
\end{equation}
Hence
\begin{align}
\ket{\psi(t)}=
\frac{1}{\sqrt{2}}\Bigg\{&
\e^{\ii Jt}\ket{\phi_-}
+\left[\cos(\Omega t)+\ii\frac{J}{\Omega}\sin(\Omega t)\right]\ket{\phi_+}
\nonumber\\
&+\ii\frac{2h}{\Omega}\sin(\Omega t)\ket{\psi_+}
\Bigg\}.
\end{align}
Returning to the computational basis gives
\begin{equation}
\boxed{
\begin{aligned}
a_{00}(t)&=\frac{1}{2}\left[\e^{\ii Jt}+\cos(\Omega t)+\ii\frac{J}{\Omega}\sin(\Omega t)\right],\\
a_{11}(t)&=\frac{1}{2}\left[-\e^{\ii Jt}+\cos(\Omega t)+\ii\frac{J}{\Omega}\sin(\Omega t)\right],\\
a_{01}(t)&=a_{10}(t)=\ii\frac{h}{\Omega}\sin(\Omega t).
\end{aligned}}
\label{eq:exactamps}
\end{equation}
The basis probabilities are $P_{ij}(t)=|a_{ij}(t)|^2$. In particular,
\begin{equation}
P_{01}(t)=P_{10}(t)=\frac{h^2}{\Omega^2}\sin^2(\Omega t).
\end{equation}

\section{Observables}

The longitudinal magnetization per spin is
\begin{equation}
M_z=\frac{Z_0+Z_1}{2}.
\end{equation}
It has eigenvalues $+1,0,0,-1$ on the ordered computational basis, so
\begin{equation}
\boxed{\langle M_z\rangle=P_{00}-P_{11}}.
\end{equation}
The nearest-neighbor spin correlation is
\begin{equation}
C_{zz}=Z_0Z_1,
\end{equation}
with eigenvalues $+1,-1,-1,+1$, hence
\begin{equation}
\boxed{\langle Z_0Z_1\rangle=P_{00}-P_{01}-P_{10}+P_{11}}.
\end{equation}
These observables compress the information in the full state. For example, $\langle M_z\rangle=0$ does not determine a unique quantum state.

\section{Why the Hamiltonian terms cannot be factorized exactly}

Split the Hamiltonian into
\begin{equation}
H=H_{ZZ}+H_X,
\end{equation}
where
\begin{equation}
H_{ZZ}=-JZ_0Z_1,\qquad H_X=-h(X_0+X_1).
\end{equation}
Using $[Z,X]=2\ii Y$,
\begin{align}
[Z_0Z_1,X_0]&=2\ii Y_0Z_1,\\
[Z_0Z_1,X_1]&=2\ii Z_0Y_1.
\end{align}
Therefore
\begin{equation}
\boxed{[H_{ZZ},H_X]=2\ii Jh\left(Y_0Z_1+Z_0Y_1\right)\neq0.}
\end{equation}
Consequently,
\begin{equation}
\e^{-\ii(H_{ZZ}+H_X)t}\neq
\e^{-\ii H_{ZZ}t}\e^{-\ii H_Xt}
\end{equation}
for a finite time interval in general.

As a state-level check,
\begin{equation}
[Z_0Z_1,X_0+X_1]\ket{10}=2\ket{00}+2\ket{11}.
\end{equation}
This is consistent with the operator identity above, but a single input-output pair does not uniquely determine the full operator.

\section{Useful exponential identities}

If an operator $P$ satisfies $P^2=\I$, then its exponential separates into even and odd powers:
\begin{align}
\e^{\ii aP}
&=\I+\ii aP+\frac{(\ii aP)^2}{2!}+\cdots\\
&=\cos(a)\I+\ii\sin(a)P.
\end{align}
This applies to every Pauli operator and to $Z_0Z_1$. Thus
\begin{equation}
\e^{\ii aX}=\cos(a)\I+\ii\sin(a)X,
\end{equation}
\begin{equation}
\e^{\ii aZ_0Z_1}=\cos(a)\I+\ii\sin(a)Z_0Z_1.
\end{equation}
Since $Z_0Z_1$ has eigenvalues $+1,-1,-1,+1$,
\begin{equation}
\e^{\ii aZ_0Z_1}=\operatorname{diag}(\e^{\ii a},\e^{-\ii a},\e^{-\ii a},\e^{\ii a}).
\end{equation}
The minus sign of an eigenvalue enters the exponent; it is not an overall prefactor.

\section{First-order Lie-Trotter formula}

For a total evolution time $t$ divided into $r$ steps,
\begin{equation}
\Delta t=\frac{t}{r}.
\end{equation}
The first-order approximation used here is
\begin{equation}
\boxed{U_1(t;r)=\left[\e^{-\ii H_{ZZ}\Delta t}\e^{-\ii H_X\Delta t}\right]^r.}
\end{equation}
A Taylor expansion shows that the exact and factorized short-time evolutions agree through first order in $\Delta t$; their difference begins at $\Order(\Delta t^2)$ and is controlled by the commutator. Repeating $r$ steps over fixed total time gives a typical global operator/state error of order
\begin{equation}
\Order\left(\frac{t^2}{r}\right).
\end{equation}

\subsection{Gate mapping}

Because $[X_0,X_1]=0$,
\begin{equation}
\e^{-\ii H_X\Delta t}
=\e^{\ii h\Delta t X_0}\e^{\ii h\Delta t X_1}.
\end{equation}
Qiskit defines
\begin{equation}
RX(\theta)=\e^{-\ii\theta X/2},
\end{equation}
so each transverse-field rotation uses
\begin{equation}
\boxed{\theta_X=-2h\Delta t.}
\end{equation}
Similarly,
\begin{equation}
\e^{-\ii H_{ZZ}\Delta t}=\e^{\ii J\Delta t Z_0Z_1},
\end{equation}
while
\begin{equation}
RZZ(\theta)=\e^{-\ii\theta Z_0Z_1/2},
\end{equation}
which gives
\begin{equation}
\boxed{\theta_{ZZ}=-2J\Delta t.}
\end{equation}

\subsection{One first-order step acting on \texorpdfstring{$\ket{00}$}{|00>}}

Define
\begin{equation}
c=\cos(h\Delta t),\qquad s=\sin(h\Delta t),
\end{equation}
\begin{equation}
p=\e^{\ii J\Delta t},\qquad q=\e^{-\ii J\Delta t}.
\end{equation}
The transverse-field part gives
\begin{align}
\e^{-\ii H_X\Delta t}\ket{00}
&=(c\ket{0}+\ii s\ket{1})\otimes(c\ket{0}+\ii s\ket{1})\\
&=c^2\ket{00}+\ii cs\ket{01}+\ii cs\ket{10}-s^2\ket{11}.
\end{align}
The $ZZ$ evolution then adds phases according to parity:
\begin{equation}
\boxed{
U_{1,\mathrm{step}}\ket{00}
=pc^2\ket{00}+\ii qcs\ket{01}+\ii qcs\ket{10}-ps^2\ket{11}.}
\label{eq:firststep}
\end{equation}
The $ZZ$ gate does not change probabilities immediately, but its relative phases affect subsequent interference after later noncommuting rotations.

\section{Second-order symmetric Suzuki-Trotter formula}

The symmetric second-order step is
\begin{equation}
\boxed{S_2(\Delta t)=
\e^{-\ii H_X\Delta t/2}
\e^{-\ii H_{ZZ}\Delta t}
\e^{-\ii H_X\Delta t/2}.}
\label{eq:S2}
\end{equation}
To verify its order, define $A=-\ii H_X$ and $B=-\ii H_{ZZ}$. Then
\begin{equation}
S_2(\Delta t)=\e^{A\Delta t/2}\e^{B\Delta t}\e^{A\Delta t/2}.
\end{equation}
Expanding each exponential through second order and multiplying while preserving operator order gives
\begin{align}
S_2(\Delta t)
={}&\I+(A+B)\Delta t\\
&+\frac{1}{2}\left(A^2+AB+BA+B^2\right)\Delta t^2
+\Order(\Delta t^3).
\end{align}
Since
\begin{equation}
(A+B)^2=A^2+AB+BA+B^2,
\end{equation}
this agrees with
\begin{equation}
\e^{(A+B)\Delta t}
=\I+(A+B)\Delta t+\frac{1}{2}(A+B)^2\Delta t^2+\Order(\Delta t^3).
\end{equation}
Therefore the local error is $\Order(\Delta t^3)$, and for fixed total time the typical global operator/state error is
\begin{equation}
\Order\left(\frac{t^3}{r^2}\right).
\end{equation}

\subsection{Gate mapping}

The two half-field evolutions use
\begin{equation}
\boxed{\theta_{X,\mathrm{half}}=-h\Delta t,}
\end{equation}
whereas the central interaction uses
\begin{equation}
\boxed{\theta_{ZZ}=-2J\Delta t.}
\end{equation}
One symmetric step is therefore represented chronologically in a circuit by
\begin{equation}
RX_0(-h\Delta t)\,RX_1(-h\Delta t)\,RZZ(-2J\Delta t)\,
RX_0(-h\Delta t)\,RX_1(-h\Delta t).
\end{equation}

\subsection{One second-order step acting on \texorpdfstring{$\ket{00}$}{|00>}}

Let
\begin{equation}
c=\cos\left(\frac{h\Delta t}{2}\right),\qquad
s=\sin\left(\frac{h\Delta t}{2}\right),
\end{equation}
and again define $p=\e^{\ii J\Delta t}$ and $q=\e^{-\ii J\Delta t}$. After the first half-field step,
\begin{equation}
\ket{\chi_1}=c^2\ket{00}+\ii cs\ket{01}+\ii cs\ket{10}-s^2\ket{11}.
\end{equation}
After the full $ZZ$ step,
\begin{equation}
\ket{\chi_2}=pc^2\ket{00}+\ii qcs\ket{01}+\ii qcs\ket{10}-ps^2\ket{11}.
\end{equation}
Applying the second half-field step and collecting equal basis states gives
\begin{align}
a_{00}&=p(c^4+s^4)-2qc^2s^2,\\
a_{01}=a_{10}&=\ii(p+q)cs(c^2-s^2),\\
a_{11}&=-2(p+q)c^2s^2.
\end{align}
Use
\begin{equation}
p+q=2\cos(J\Delta t),
\end{equation}
\begin{equation}
2cs=\sin(h\Delta t),\qquad c^2-s^2=\cos(h\Delta t),
\end{equation}
and
\begin{equation}
4c^2s^2=\sin^2(h\Delta t).
\end{equation}
The compact result is
\begin{equation}
\boxed{
\begin{aligned}
a_{00}&=\cos(J\Delta t)\cos^2(h\Delta t)+\ii\sin(J\Delta t),\\
a_{01}=a_{10}&=\ii\cos(J\Delta t)\sin(h\Delta t)\cos(h\Delta t),\\
a_{11}&=-\cos(J\Delta t)\sin^2(h\Delta t).
\end{aligned}}
\label{eq:secondstep}
\end{equation}
The equality $a_{01}=a_{10}$ follows from exchange symmetry and is also a useful debugging condition.

\subsection{Limiting-case checks}

If $h=0$, Eq.~\eqref{eq:secondstep} becomes
\begin{equation}
\ket{\psi}=\e^{\ii J\Delta t}\ket{00}.
\end{equation}
If $J=0$, it becomes the exact product of two independent $X$ rotations:
\begin{align}
\ket{\psi}={}&\cos^2(h\Delta t)\ket{00}
+\ii\sin(h\Delta t)\cos(h\Delta t)(\ket{01}+\ket{10})\\
&-\sin^2(h\Delta t)\ket{11}.
\end{align}
The amplitudes also satisfy
\begin{equation}
|a_{00}|^2+|a_{01}|^2+|a_{10}|^2+|a_{11}|^2=1.
\end{equation}

\section{Error measures and convergence}

For normalized pure states, the fidelity is
\begin{equation}
F(\psi_{\mathrm{exact}},\psi_{\mathrm{approx}})
=\left|\braket{\psi_{\mathrm{exact}}}{\psi_{\mathrm{approx}}}\right|^2.
\end{equation}
The infidelity is $1-F$. A global phase leaves the fidelity equal to one even though direct component-wise comparison may fail.

For a first-order product formula, the typical state/operator error scales as $r^{-1}$; for the symmetric second-order formula it scales as $r^{-2}$. In a generic small-error regime, the infidelity is quadratic in the perpendicular state error, leading commonly to
\begin{equation}
1-F_1\sim r^{-2},\qquad 1-F_2\sim r^{-4}.
\end{equation}
These are asymptotic expectations rather than guarantees at every coarse step size or special evolution time. Coherent cancellation can make finite-step behavior non-monotonic.

Observable errors are defined separately, for example
\begin{equation}
\Delta M_z=
\left|\langle M_z\rangle_{\mathrm{approx}}-
\langle M_z\rangle_{\mathrm{exact}}\right|.
\end{equation}
A small error in one observable does not imply that the complete statevector is accurate.

\section{Circuit cost}

A first-order step contains two $RX$ gates and one $RZZ$ gate. An explicit second-order step contains four $RX$ gates and one $RZZ$ gate. Therefore, comparing methods at equal step count is not the same as comparing them at equal circuit cost.

In consecutive second-order steps, adjacent half-field rotations can be combined:
\begin{equation}
RX(-h\Delta t)RX(-h\Delta t)=RX(-2h\Delta t).
\end{equation}
Before hardware execution, further decomposition and transpilation would be required, and the cost of $RZZ$ depends on the available native gate set and connectivity.

\section{Numerical reference point}

For the values used frequently in the notebooks,
\begin{equation}
J=1.0,\qquad h=0.7,
\end{equation}
we find
\begin{equation}
\Omega=\sqrt{1+4(0.7)^2}\approx1.72046505,
\end{equation}
and the exact spectrum is approximately
\begin{equation}
(-1.72046505,-1,1,1.72046505).
\end{equation}
For one symmetric second-order step with $\Delta t=1$,
\begin{equation}
\ket{\psi_{S_2}}\approx
(0.31606797+0.84147098\ii)\ket{00}
+0.26622038\ii\ket{01}
+0.26622038\ii\ket{10}
-0.22423433\ket{11}.
\end{equation}
This reference value is useful for checking parameter ordering and gate-angle conventions.

\section{Implementation checklist}

A reliable notebook workflow should include the following independent checks:
\begin{enumerate}[leftmargin=1.5em]
\item Verify the action of each tensor-product operator on all four basis states.
\item Confirm that $H=H^\dagger$.
\item Compare numerical eigenvalues with $(-\Omega,-J,J,\Omega)$.
\item Compare exact matrix evolution with Eq.~\eqref{eq:exactamps}.
\item Compare one-step Qiskit statevectors with Eqs.~\eqref{eq:firststep} and \eqref{eq:secondstep}.
\item Check normalization and exchange symmetry $a_{01}=a_{10}$.
\item Print the actual $RX$ and $RZZ$ angles when debugging.
\item Distinguish a true component mismatch from an irrelevant global phase by computing fidelity.
\item Restart the notebook kernel and run all cells before committing results.
\end{enumerate}

\section{Summary}

The two-qubit model is small enough to solve exactly but already contains the central ingredients of digital quantum simulation: noncommuting Hamiltonian terms, tensor-product operators, basis-dependent measurements, product-formula errors, and a trade-off between accuracy and circuit resources. The exact solution provides a controlled reference against which the gate-level approximations can be validated. The first-order formula supplies the simplest digital decomposition, while the symmetric second-order formula demonstrates how time-symmetric operator ordering cancels the leading local error.

\end{document}
